% ==============================================================================
% Chapter 1: Introduction
% ==============================================================================
\chapter{Introduction}
\label{ch:introduction}

\section{Aim and Motivation}
\label{sec:aim_motivation}

Informal street vending constitutes a fundamental socioeconomic pillar of urban economies throughout the Global South, providing self-employment, micro-entrepreneurial livelihood, and daily sustenance to millions of vulnerable urban residents~\cite{bhowmik2005street}. In fast-growing Indian tier-II metropolitan centres such as Nashik, Maharashtra, street vendors supply affordable daily necessities---including fresh produce, street food, garments, and household staples---directly to consumers within dense residential and commercial neighborhoods.

However, the rapid expansion of informal vending in unplanned urban spaces precipitates acute friction with municipal urban infrastructure~\cite{bromley2000street}. Unregulated clustering of vendors around busy transit intersections, railway stations, and arterial road networks results in severe pedestrian sidewalk encroachment, forced vehicular spillover, public right-of-way bottlenecks, and vehicular congestion. Historically, municipal corporations have responded to these challenges through cyclical, heavy-handed eviction drives. These coercive enforcement methods fail to resolve congestion permanently; evicted vendors invariably return within hours due to lack of viable alternative locations and acute economic necessity.

Recognizing the indispensable economic contribution of street vendors and the systemic failure of punitive evictions, the Parliament of India enacted statutory legislation: the \textit{Street Vendors (Protection of Livelihood and Regulation of Street Vending) Act, 2014}~\cite{india2014street}. The Act mandates that urban local bodies (ULBs):
\begin{enumerate}
    \item Conduct periodic biometric surveys to register and issue vending certificates to all informal vendors;
    \item Demarcate urban jurisdictions into designated \textbf{Free Vending Zones}, \textbf{Restricted Vending Zones}, and strictly demarcated \textbf{No-Vending Zones};
    \item Formulate Town Vending Committees (TVCs) to establish statutory zone carrying capacities ($C_z$);
    \item Guarantee that no street vendor is arbitrarily evicted or relocated until suitable alternative accommodation within an authorized vending zone is identified and allotted.
\end{enumerate}

Despite this statutory mandate, municipal authorities across India---including the Nashik Municipal Corporation (NMC)---face a severe computational and operational void. Town-planning departments continue to manage vendor records using fragmented spreadsheets, isolated PDF municipal gazettes, and subjective manual allocation processes. Crucially, town-planning officers lack objective, data-driven geospatial tools to:
\begin{itemize}
    \item Quantitatively evaluate whether a designated municipal zone provides adequate economic footfall, pedestrian accessibility, and public utility infrastructure for specific vending categories;
    \item Systematically recommend equitable relocation options when high-density zones undergo civil redevelopment or anti-encroachment restructuring;
    \item Enforce hard capacity constraints during batch relocation to prevent algorithmic over-saturation of popular commercial centres;
    \item Demarcate metric, physical micro-slots on pedestrian walkways that satisfy urban safety clearance corridors and prevent obstruction of building frontages.
\end{itemize}

The central aim of \textbf{VendorMap} (UrbanVending AI) is to bridge this critical gap between statutory urban policy and operational civic reality. By fusing Geographic Information Systems (GIS), Multi-Criteria Decision Analysis (MCDA), mathematical Mixed-Integer Linear Programming (MILP), and full-stack software engineering, VendorMap delivers an automated, transparent, and legally compliant decision-support platform for municipal vendor management, equitable relocation, and physical micro-slot demarcation.

\section{Scope of the Project}
\label{sec:scope}

The operational and research scope of VendorMap encompasses the complete administrative boundary of the Nashik Municipal Corporation (NMC), encompassing six administrative divisions: \textit{Nashik East, Nashik West, Nashik Road, Panchavati, CIDCO, and Satpur}. The project specifically addresses:

\begin{enumerate}
    \item \textbf{Municipal Registry Ingestion and Data Normalization:} Curating and harmonizing official civic datasets comprising 7,259 registered street vendors and 308 municipal vending zones (categorized into 166 FREE, 59 RESTRICTED, and 83 NO\_VENDING zones).
    \item \textbf{OpenStreetMap GIS Spatial Feature Extraction:} Engineering 12 multi-scale environmental indicators across 13,988 road vectors and 353 commercial/transit Points-of-Interest (POIs) using logarithmic saturation and exponential proximity decay transformations.
    \item \textbf{Spatial Environmental Typology:} Employing unsupervised $K$-Means clustering ($K=2$) to establish reproducible macro-typologies across 258 usable zone geometries without confounding statistical clusters with statutory zoning bylaws.
    \item \textbf{Multi-Criteria Decision Analysis (MCDA):} Developing an operational six-factor decision engine integrating commercial density, accessibility, public amenities, vendor category compatibility, spatial preference matching, and historical vendor concentration.
    \item \textbf{Empirical Bayes Category Shrinkage:} Implementing leave-location-group-out cross-validated Laplace shrinkage ($M=20.0$) to prevent sample-size distortion and neutralize extreme lift values for rare vendor categories in small zones.
    \item \textbf{Analytic Hierarchy Process (AHP) Policy Calibration:} Providing town-planning officers with an interactive pairwise comparison interface with real-time eigenvalue decomposition and Consistency Ratio ($CR \le 0.10$) validation.
    \item \textbf{Pre-emptive Goal Programming Optimizer:} Formulating a two-stage Mixed-Integer Linear Program (MILP) solved via the HiGHS dual simplex solver to eliminate the \textit{Spatial Capacity Paradox} by maximizing placed vendors while strictly respecting municipal zone capacity ceilings.
    \item \textbf{Metric Micro-Slot Packing Generator:} Developing a sub-millimeter geometric packing engine in UTM Zone 43N projection (EPSG:32643) that generates $2\,\text{m} \times 2\,\text{m}$ vending stalls with $\ge 0.5\,\text{m}$ inter-slot spacing and an uninterrupted $1.2\,\text{m}$ pedestrian clearance corridor, empirically validated on the Lekha Nagar pilot site (Zone NN\_F\_013).
    \item \textbf{Full-Stack Operational Delivery Platform:} Engineering an authenticated React 18 / FastAPI web application with role-based access control, SQLite Write-Ahead Logging, and in-app turn-by-turn walking navigation.
    \item \textbf{Privacy-Preserving Vendor Identity Protection:} Eliminating plaintext phone number storage by enforcing HMAC-SHA256 one-way cryptographic hashing keyed by institutional secrets.
    \item \textbf{Direct Meta WhatsApp Cloud API Integration:} Providing a low-barrier, bilingual (English and Marathi) conversational channel for vendor onboarding, OTP verification, and status inquiry without third-party wrapper dependencies.
\end{enumerate}

\section{Objectives}
\label{sec:objectives}

The specific engineering, research, and implementation objectives of the VendorMap project are:

\begin{enumerate}
    \item \textbf{Objective 1 (Data Integration \& Geospatial Feature Pipeline):} To ingest, clean, and spatialize the official NMC registry of 7,259 vendors and 308 zones, extracting 12 non-linear OpenStreetMap GIS environmental indicators spanning commercial footfall, transit proximity, and public facility infrastructure.
    \item \textbf{Objective 2 (Environmental Typology Modeling):} To implement unsupervised $K$-Means clustering across usable municipal zones and establish cluster partition stability across multiple random seeds, achieving an Adjusted Rand Index of $ARI \ge 0.95$.
    \item \textbf{Objective 3 (AHP-Calibrated Multi-Criteria Decision Engine):} To develop a missing-data-aware six-factor MCDA recommendation engine with leave-group-out empirical Bayes shrinkage ($M=20.0$), calibrated via AHP with $CR \le 0.10$, and verify ranking robustness across an extensive sensitivity audit ($\ge 200$ scenarios).
    \item \textbf{Objective 4 (Constrained Optimization \& Paradox Resolution):} To formulate and implement a two-stage pre-emptive Goal Programming optimizer (MILP) that resolves the Spatial Capacity Paradox, guaranteeing 0\% capacity violation while maximizing total accommodated vendors.
    \item \textbf{Objective 5 (Sub-Millimeter Micro-Slot Packing):} To design a geometric micro-allocation algorithm in projected coordinates (EPSG:32643) that packs standardized $2\,\text{m} \times 2\,\text{m}$ pitches within arbitrary municipal boundary polygons, preserving $\ge 0.5\,\text{m}$ clearance and a continuous $1.2\,\text{m}$ pedestrian corridor with edge deviations $< 0.002\,\text{m}$.
    \item \textbf{Objective 6 (Operational Full-Stack Architecture):} To build a production-ready, authenticated web platform using FastAPI and React 18 with SQLite ACID transactions, providing dedicated administrative dashboards for municipal officers and self-service portals for registered vendors.
    \item \textbf{Objective 7 (Cryptographic Data Minimization):} To design and enforce an end-to-end cryptographic privacy architecture utilizing HMAC-SHA256 identity hashing, ensuring that sensitive vendor contact details are never stored or exposed in plaintext.
    \item \textbf{Objective 8 (Conversational Multilingual Channel):} To integrate Meta's official WhatsApp Cloud API directly with an asynchronous state machine supporting English and Marathi dialogue for vendor self-service and status tracking.
    \item \textbf{Objective 9 (In-App Turn-by-Turn Navigation):} To implement single-map vendor walking guidance utilizing browser GPS geolocation and open-source routing providers, directing approved vendors directly to their designated micro-slot coordinates.
\end{enumerate}

\section{Report Organization}
\label{sec:report_organization}

The remainder of this Project Stage-II Report is organized into the following eight chapters:

\begin{itemize}
    \item \textbf{Chapter 2 (Literature Review):} Reviews foundational academic literature on street vending governance, Multi-Criteria Decision Analysis in urban planning, the Analytic Hierarchy Process, spatial clustering, and facility location optimization. Critically evaluates existing municipal practices and identifies technical research gaps.
    \item \textbf{Chapter 3 (System Architecture):} Formulates the formal mathematical problem statement, presents the end-to-end seven-stage methodology, details the comprehensive multi-tier system block diagram, and provides algorithmic pseudocode for the complete allocation pipeline.
    \item \textbf{Chapter 4 (Project Requirement Specifications):} Outlines the hardware and software development and operational requirements, enumerates functional and non-functional requirements, and states core system constraints and domain assumptions.
    \item \textbf{Chapter 5 (High-Level Design):} Presents detailed architectural models including the relational Entity-Relationship Diagram (ERD), Level 0 and Level 1 Data Flow Diagrams (DFDs), and UML diagrams (Use Case, Service Class Architecture, Sequence, and Activity diagrams).
    \item \textbf{Chapter 6 (Experimental Setup):} Details the official NMC and OpenStreetMap datasets, geometric preprocessing workflows, experimental evaluation protocols, quantitative performance parameters, and computational efficiency considerations.
    \item \textbf{Chapter 7 (Results and Evaluation):} Reports quantitative findings across spatial clustering, 224-scenario MCDA sensitivity analysis, Goal Programming benchmarks, Lekha Nagar micro-slot geometric validation, supervised regression models with TreeSHAP explainability, automated test case results, GUI module walkthroughs, feature evaluation, and an operational cost analysis.
    \item \textbf{Chapter 8 (Project Planning):} Details the 16-week Stage-II project execution plan, weekly activity schedule, project milestones, and a comprehensive Gantt chart.
    \item \textbf{Chapter 9 (Conclusions):} Summarizes the technical contributions and policy implications of the implemented system, identifies current operational limitations, and outlines future research directions.
    \item \textbf{Back Matter:} Contains the IEEE-formatted Bibliography and Appendices spanning plagiarism report placeholders, base paper analyses, tool inventories, and publication records.
\end{itemize}
